\section{What does not explain the split}
\label{sec:controls}

\begin{figure*}[t]
\centering
\includegraphics[width=\textwidth]{fig2_controls}
\caption{Three alternative explanations, each held fixed. (a) The internal advantage on the BFCL checkpoints, raw and recomputed within strata of item difficulty, where difficulty is the share of the other checkpoints that fail the same item. (b) The token-role probe's AUC on the two Llama-3.2 runs as its training failures are thinned to the counts the other families supply, with the mean log-probability dashed. (c) The advantage against the mean log-probability and against the confidence summary chosen on training folds, one point per run with stored summaries.}
\label{fig:controls}
\end{figure*}

A difference between model families is the kind of result that is usually something else, so each explanation that could produce it from the stored data was tested before it was believed. One of them, label error, did produce part of it and is reported first.

\paragraph{Labels and extraction.}
Before the audit of \cref{sec:setup}, Qwen3-1.7B appeared on the side where internals are needed, with an advantage of \gapQwenThreeBeforeAudit{}. Of its \repPosBeforeQwenThreeBfcl{} labelled failures on the scored population, \repPosAfterQwenThreeBfcl{} survive the repair of the list-argument comparison, and its advantage falls to \GapQwenThreeBfcl{}. The defect could not have produced the other side: it changed at most \labelChangedRecentMax{} labels on any Qwen3.5 or MiniCPM5 run, and the probe on MiniCPM5, which had read the end-of-sequence state for all three roles, moves by two points once it reads the call. On the re-extracted runs the probe's positions are located on \locatedMin{}\% or more of scored items, and the token-level probe, which needs none, orders the sides the same way (\cref{sec:result}). A reviewer should read the earlier family effect as an artefact and the present one as what survives its removal.

\paragraph{Item difficulty.}
A model that fails rarely fails on hard items, and its confidence may register that an item is hard rather than that its own call is wrong. The items of BFCL are the same text for every model, so an item's difficulty is the share of the other checkpoints that fail it, and that share alone predicts a model's failures at \diffOnlyMin{} to \diffOnlyMax{} AUC. Recomputing the probe and the confidence inside strata of that difficulty, weighted by the positive-negative pairs each stratum holds, removes everything either detector gains from knowing which items are hard. The split widens (\cref{fig:controls}a): the advantage moves from \diffGapRawLlamaOneB{} to \diffGapWithinLlamaOneB{} on Llama-3.2-1B and from \diffGapRawMiniCpm{} to \diffGapWithinMiniCpm{} on MiniCPM5-2B. Confidence on the models where it suffices is tracking the model's own errors, not the hardness of the request.

\paragraph{Label budget.}
A probe needs labelled failures and confidence needs none, and the models where confidence suffices fail less often, supplying \nPosQwenThreeBfcl{} to \nPosQwenThreeFiveBfcl{} failures against hundreds on Llama. Training the Llama probes on exactly those counts, with the test folds untouched, costs them \budgetProbeDropLlamaOneB{} AUC at 1B and \budgetProbeDropLlamaThreeB{} at 3B (\cref{fig:controls}b). At \budgetMinPosLlamaOneB{} training failures the advantage is still \budgetGapAtMinLlamaOneB{} and \budgetGapAtMinLlamaThreeB{}, both above every value on the other side at any budget.

\paragraph{The confidence summary.}
The mean log-probability is one summary of the output distribution. On the runs that store eleven length-invariant summaries, the one with the highest AUC on the other folds' items is applied fold by fold, so no summary is chosen on the items it is scored on (\cref{fig:controls}c). Two runs on the side where internals are needed store several summaries: on Llama-3.2-1B on BFCL the fold-chosen summary moves the advantage from \GapLlamaOneBBfcl{} to \gapBestLlamaOneBBfcl{}, and on Gemma-3-1B on BFCL the fold-chosen summary generalises worse than the mean and the advantage rises to \gapBestGemmaBfcl{}. On the other side the largest value is \gapBestConfidenceMax{}. An oracle that picks the best single summary on the test items, which no deployment could do, leaves the first side at \gapOracleMinInternals{} to \gapOracleMaxInternals{} and the second at \gapOracleMinConfidence{} to \gapOracleMaxConfidence{}. Two summaries that grow with the length of the call, its token count and its total log-probability, are excluded from the candidates because they are length features, which the floor already covers.

\paragraph{Benchmark.}
Gemma-3 reads \GapGemmaGlaive{} on Glaive and \GapGemmaBfcl{} on BFCL, Llama-3.2-1B reads \GapLlamaOneBGlaive{}, \GapLlamaOneBBfcl{} and \GapLlamaOneBLive{} across three benchmarks, and MiniCPM5 reads \GapMiniCpmBfcl{} and \GapMiniCpmLive{} on two, so no run's side depends on its benchmark. Qwen3 and MiniCPM5 on Glaive are underpowered, with \nPosQwenThreeGlaive{} and \nPosMiniCpmGlaive{} failures in several hundred calls, and their intervals span zero. That they almost never fail there is itself consistent with their side.

\paragraph{Training population and pooling.}
Two protocol choices move the numbers and are fixed as stated in \cref{sec:setup}. Training the probe on every item while scoring the call-required population teaches it the format failures that population excludes, and on BFCL the same behaviour, answering without a call, is a failure on one category and a success on another. Training on the scored population removes a spurious family difference on the population that includes abstentions and changes the call-required numbers by less than the replication drift. Fixing the sign of the log-probability on the scored population rather than on every item raises confidence on two Llama runs by up to a quarter of the AUC scale without moving their side. Pooling out-of-fold scores across folds with different calibrations penalises a fitted probe and not a fixed confidence score: the advantage computed within folds is higher than the pooled one by \poolShiftMean{} on average and \poolShiftMax{} at most, and ranges from \gapWithinMinInternals{} to \gapWithinMaxInternals{} on the first side and \gapWithinMinConfidence{} to \gapWithinMaxConfidence{} on the second.

\paragraph{A registered prediction that failed.}
Before the audit, the split had been written down as a division by release period, the two 2026 families against the three earlier ones, with a decision rule committed to the repository. That prediction failed: Qwen3, released in April 2025, sits with the 2026 families once the labels are right. A cutoff placed one month later would have passed, which is why the division that holds is reported as a description and the property behind it as a hypothesis (\cref{app:registry}).
